\documentclass[12pt]{article}

% ============================================================================
% AER-Style Preamble
% ============================================================================
\Sepackage{fontspec}
\
\setmainfont{Times New Roman}Sepackage{unicode-math}setmathfont{XITS Math}
\usepackage{amsmath,amssymb,amsthm}
\usepackage{graphicx}
\usepackage{booktabs}
\usepackage{natbib}
\usepackage{setspace}
\usepackage[margin=1in]{geometry}
\usepackage{caption}
\usepackage{subcaption}
\usepackage{float}
\usepackage{hyperref}
\usepackage{xcolor}
\usepackage{multirow}
\usepackage{array}
\usepackage{threeparttable}
\usepackage{rotating}
\usepackage{pdflscape}
\usepackage{enumitem}

\hypersetup{colorlinks=true, linkcolor=blue!60!black, citecolor=blue!60!black, urlcolor=blue!60!black}
\captionsetup{font=small, labelfont=bf, labelsep=period}

\doublespacing
\setlength{\parindent}{1.5em}
\setlength{\parskip}{0pt}

\newtheorem{proposition}{Proposition}
\newtheorem{corollary}{Corollary}

% ============================================================================
% Title Page
% ============================================================================

\title{\textbf{The Cost of Closing the Strait of Hormuz: A DSGE Analysis}\thanks{We thank the editors and two anonymous referees for valuable comments. This paper builds on the empirical framework of \citet{federle2026price}. All errors are our own.}}

\author{
  \textsc{Wang Weijia} (王维佳)\thanks{School of Economics and Finance, Xi'an Jiaotong University, Xi'an, Shaanxi 710000, China. Email: wangweijia@stu.xjtu.edu.cn. Wang Weijia (1993.08--), from Baoding, Hebei Province, Ph.D.\ candidate at the School of Economics and Finance, Xi'an Jiaotong University. Research interests: international trade.}
}

\date{April 2026}

\begin{document}

\maketitle

\begin{abstract}
\noindent
How large are the global welfare costs of a Strait of Hormuz closure, and which transmission channels drive them? We build a four-region New Keynesian DSGE model of a US--Israel conflict with Iran, integrating trade, energy, financial, and stabilization channels, and validate the financial linkage matrix with a structural VAR(1) estimated on four historical Middle East conflict episodes. A complete Hormuz blockade collapses Iran's GDP by 27.2 percent (34.8 percent under second-order perturbation), spikes global oil prices by 63 percent, and inflicts 4.88 percentage points of global welfare loss through financial frictions alone---exceeding the entire trade-channel effect; conventional monetary policy moves global welfare by less than 0.11 percentage points across all realistic Taylor-rule regimes, confirming its impotence against rare-disaster supply shocks; and China's manufacturing capacity, not its strategic petroleum reserve, is the binding stabilization margin, with reserve releases hitting their physical ceiling by quarter one while doubling the manufacturing response saves 0.70 percentage points of China's own welfare. Supply-chain resilience and macroprudential policy, not central-bank activism, are the correct policy levers for geopolitical rare disasters.

\bigskip
\noindent \textbf{JEL Codes:} E32, F41, F51, H56, Q43 \\
\noindent \textbf{Keywords:} Strait of Hormuz, DSGE, oil shocks, financial contagion, China stabilization, rare disasters, monetary policy limits
\end{abstract}

\thispagestyle{empty}
\newpage
\setcounter{page}{1}


% ============================================================================
% I. INTRODUCTION
% ============================================================================

\section{Introduction}

Closing the Strait of Hormuz would cut off 20 percent of global petroleum trade in a single chokepoint. Yet the macroeconomics literature has no structural model of what such a closure would actually cost---or through which channels those costs would flow. \citet{federle2026price} document that wars reduce war-site GDP by roughly 10 percent and generate third-country spillovers through bilateral trade, but their reduced-form local projections cannot trace the energy, financial, and stabilization channels that activate when a single chokepoint seizes global commodity flows, nor can they answer the policy counterfactuals that motivate current strategic planning: How much does financial contagion amplify trade losses? Does aggressive monetary easing help? Would a Chinese stabilization effort matter, and if so, through which margin?

We answer these questions with a four-region New Keynesian DSGE model calibrated to a US--Israel military strike on Iran that triggers a Hormuz blockade. The four regions---Iran (war site), the US--Israel bloc (belligerent), China (stabilizer), and the rest of the world (ROW)---are each governed by a New Keynesian three-equation system and linked through a global oil market, bilateral trade flows, and a financial contagion matrix. We validate the financial linkage matrix empirically: a structural VAR(1) estimated on quarterly sovereign-spread data from four historical Middle East conflict episodes recovers transmission coefficients of $+0.57$ (Iran to US, $t = 11.4$) and $+0.50$ (Iran to ROW, $t = 5.8$), tightly bracketing our hand-calibrated values and confirming that the financial channel is real, not an artifact of parameterization. The model is then re-solved under second-order perturbation to quantify linearization bias, and nine policy counterfactuals map the welfare frontier across monetary, financial, and stabilization regimes.

\paragraph{Three contributions.} Our structural analysis delivers three quantitative results that sharpen the policy debate in ways reduced-form estimates cannot.

\textit{First, financial frictions---not goods trade---are the dominant amplifier of a Hormuz closure.} Sovereign risk premia, capital-flow reversals, and equity-market contagion raise ROW peak output losses from 1.1 percent (trade channel alone) to 3.8 percent, and account for 4.88 percentage points of cumulative global consumption-equivalent welfare loss---more than the entire trade channel. The VAR evidence confirms this is not a calibration choice: geopolitical shocks spread primarily through risk-off portfolio flows, as \citet{caldara2022measuring} document empirically and as our structural framework now embeds with full welfare accounting. Macroprudential policy that limits cross-border contagion would therefore save more global welfare than any conceivable monetary intervention.

\textit{Second, conventional monetary policy is impotent against rare-disaster supply shocks.} Nine counterfactuals spanning $\phi_\pi \in [1.0, 2.5]$, coordinated demand support, and selective financial-channel suppression reveal that the global welfare difference between an aggressive and a dovish central bank is less than 0.11 percentage points. A 27 percent war-site GDP collapse and a 63 percent oil price spike together exhaust the demand-management capacity of any standard Taylor rule. The result sharpens the rare-disaster literature \citep{barro2006rare}: in catastrophic supply shocks, the policy lever is not the interest rate but supply-chain resilience and energy security.

\textit{Third, China's manufacturing capacity---not its strategic petroleum reserve---is the binding stabilization margin.} China's triple mechanism (SPR releases, renewable substitution, manufacturing expansion) reduces cumulative global GDP losses by \$2.2 trillion and cuts peak global inflation by 0.29 percentage points. Under endogenous reaction functions, the SPR rule hits its physical capacity ceiling from the first quarter onward, making additional reserves worthless at the margin. Doubling the manufacturing reaction coefficient saves 0.70 percentage points of China's own welfare. This finding extends the ``China shock'' literature \citep{autor2013china} in an unexpected direction: in geopolitical crises, China's diversified industrial base functions as a global supply-chain buffer, generating positive externalities for third countries rather than competitive disruption.

\paragraph{Robustness.} The 5.9 percentage-point linearization bias we detect via second-order perturbation is concentrated in the war site and does not alter the cross-regional spillover pattern, the monetary-policy impotence result, or the manufacturing-over-reserves ranking. All three findings survive replacement of the calibrated financial matrix with the VAR-estimated matrix (welfare numbers shift by less than 5 percent) and replacement of deterministic stabilization paths with endogenous reaction functions. Section~\ref{sec:robustness} documents these checks; the Online Appendix provides the full VAR estimation and perturbation algorithm.

Sections~\ref{sec:literature}--\ref{sec:calibration} present the literature, model, and calibration. Section~\ref{sec:results} reports baseline results. Section~\ref{sec:robustness} presents robustness exercises. Section~\ref{sec:conclusion} draws policy implications.


% ============================================================================
% II. LITERATURE REVIEW
% ============================================================================

\section{Related Literature}
\label{sec:literature}

Four strands of the literature bear directly on our analysis.

\paragraph{The economics of war and conflict.} Wars reduce war-site GDP by roughly 10 percent on average; the spillover to third countries flows primarily through bilateral trade exposure \citep{federle2026price}. Reduced-form local projections identify these magnitudes cleanly but cannot decompose them into energy, financial, and supply-chain components, nor can they counterfactually vary policy. \citet{abadie2003economic} and \citet{barro2006rare} establish the foundational methodology---synthetic control and rare-disaster pricing---that we extend into a multi-channel structural framework. Our model uses \citeauthor{federle2026price}'s estimated moments as calibration targets and embeds them in a DSGE environment that permits the counterfactual analysis their reduced-form specification rules out by design.

\paragraph{DSGE models of geopolitical risk.} \citet{caldara2022measuring} show that geopolitical risk shocks reduce investment and output significantly, and that the transmission is concentrated in financial markets rather than goods trade---the same channel ordering our model confirms. \citet{bodenstein2011oil} develop the two-country open-economy DSGE framework with oil that we extend to four regions; their calibration of short-run oil elasticities ($|\varepsilon_d| = 0.05$, $\varepsilon_s = 0.10$) is standard in the literature and carries directly into our oil-price equation~(\ref{eq:oil_price}). The key innovation in our paper relative to both is an empirically estimated cross-regional financial contagion matrix, which we validate using VAR evidence rather than assuming.

\paragraph{Oil shocks and the macroeconomy.} \citet{hamilton1983oil} establishes that oil supply disruptions cause recessions; \citet{kilian2009not} shows the mechanism depends on whether the shock is supply- or demand-driven; and \citet{blanchard2007macroeconomic} document that reduced energy intensity and better monetary policy have dampened oil-shock pass-through since the 1980s. The Hormuz scenario is a textbook supply disruption---exogenous to demand, concentrated in a chokepoint, and larger in magnitude than any peacetime disruption---making it the cleanest possible test case for these transmission channels. Our model inherits \citeauthor{blanchard2007macroeconomic}'s finding that energy-intensity heterogeneity across regions matters for pass-through differentials, which we capture through the region-specific $\alpha_{e,i}$ parameters.

\paragraph{China's global economic role.} \citet{autor2013china} document that China's export expansion generates large negative employment effects in advanced economies. \citet{helveston2019china} show that China's solar supply chain has dramatically reduced global renewable energy costs, an externality of the same type---cost reduction for third countries---that our manufacturing stabilization channel formalizes in a crisis context. \citet{zhang2020china} analyze China's SPR using a DSGE model and find reserve releases effective at dampening oil price volatility. Our finding that SPR releases hit their physical ceiling while manufacturing capacity remains slack extends \citeauthor{zhang2020china}'s framework by modeling the intensive margin of each stabilization channel separately and ranking them by welfare impact.

\paragraph{Financial contagion and sovereign risk.} The financial channels in our model build on the literature on sovereign risk premia and capital flow dynamics. \citet{longstaff2011financial} study contagion in sovereign credit default swap markets, while \citet{rey2015dilemma} documents the global financial cycle driven by US monetary policy and risk appetite. Our financial contagion matrix captures the cross-regional transmission of sovereign spreads, with amplification through a global risk aversion factor that responds to capital flight from the conflict zone.

\paragraph{China's role in the global economy.} The literature on China's global economic impact has largely focused on the ``China shock'' in trade \citep{autor2013china} and China's role in global imbalances \citep{bernanke2005global}. More recent work examines China's strategic petroleum reserves \citep{zhang2020china} and industrial policy in renewable energy \citep{helveston2019china}. We contribute a novel perspective by modeling China not as a source of disruption but as a potential stabilizer during global crises---a role enabled by its large SPR, expanding renewable capacity, and diversified manufacturing base.


% ============================================================================
% III. MODEL
% ============================================================================

\section{The Model}
\label{sec:model}

The Strait of Hormuz closure creates three analytically distinct transmission channels: a supply-cost channel (oil prices), a demand-destruction channel (trade disruption and capital flight), and a contagion channel (sovereign spreads spreading risk across borders). Capturing all three within a unified welfare framework requires a model that is simultaneously multi-regional, open-economy, and financially augmented. We build this model as a four-region New Keynesian DSGE in which regions are linked through a global oil market, bilateral trade flows, and a financial contagion matrix. Each region $i \in \{\text{Iran}, \text{US-Israel}, \text{China}, \text{ROW}\}$ is characterized by a representative household, a continuum of monopolistically competitive firms, and a central bank following a Taylor rule. Three modeling choices distinguish the framework from a standard open-economy New Keynesian model, and each is motivated by direct empirical evidence.

\begin{enumerate}[leftmargin=2em, itemsep=2pt, topsep=4pt]
\item \textit{Chokepoint oil market.} Rather than treating oil as a symmetric bilateral trade flow, we model the Strait of Hormuz as a transit node through which 20 percent of global supply routes, so that a blockade disrupts all transit---not just Iran's bilateral exports. This is necessary because the welfare costs of the scenario operate primarily through the global oil-price channel, not Iran's bilateral trade.
\item \textit{Financial contagion matrix.} We augment the standard IS--PC--TR block with sovereign spreads, capital flows, and equity prices linked across regions by a $4 \times 4$ financial transmission matrix. This is motivated by \citet{caldara2022measuring}'s finding that geopolitical risk transmits primarily through financial markets, and is validated here by the VAR evidence in Section~\ref{sec:var}.
\item \textit{China's triple stabilization.} We model China's SPR, renewable substitution, and manufacturing expansion as three distinct policy instruments with different temporal profiles and capacity constraints. This distinction is necessary to identify which margin binds---a question that a single composite stabilization parameter cannot answer.
\end{enumerate}

\subsection{Regional Equilibrium}

Three equations govern each region: a New Keynesian Phillips curve (supply), an augmented IS curve (demand), and a Taylor rule (monetary policy). We present each equation and justify the non-standard additions.

\subsubsection{New Keynesian Phillips Curve}

Supply dynamics follow a Calvo-pricing New Keynesian Phillips curve augmented with an energy pass-through term and a physical destruction term. The energy term is necessary because the oil price spike acts as a cost-push shock whose magnitude is proportional to each region's energy intensity; omitting it would force the model to attribute oil-driven inflation to demand, reversing the identification of the monetary-policy experiment in Section~\ref{sec:counterfactuals}. The supply side of each region is governed by:
\begin{equation}
\pi_{i,t} = \beta \, \mathbb{E}_t[\pi_{i,t+1}] + \kappa_i \, \hat{y}_{i,t-1} + \xi^{\text{supply}}_{i,t}
\label{eq:nkpc}
\end{equation}
where $\pi_{i,t}$ is inflation, $\hat{y}_{i,t}$ is the output gap, $\beta$ is the discount factor, and $\kappa_i$ is the region-specific slope of the Phillips curve:
\begin{equation}
\kappa_i = \frac{(1-\theta_i)(1-\beta\theta_i)}{\theta_i} \cdot \frac{\sigma_i + \phi_i}{1 + \epsilon_i \phi_i}
\label{eq:kappa}
\end{equation}
with $\theta_i$ the Calvo price stickiness parameter, $\sigma_i$ the coefficient of relative risk aversion, $\phi_i$ the Frisch elasticity inverse, and $\epsilon_i$ the elasticity of substitution across varieties. The supply shock $\xi^{\text{supply}}_{i,t}$ combines energy cost pass-through and physical destruction:
\begin{equation}
\xi^{\text{supply}}_{i,t} = \alpha_{e,i} \cdot \max(\Delta p^{\text{oil}}_t, 0) \cdot 0.6 + \delta^{\text{destroy}}_{i,t} \cdot 0.8 - \chi^{\text{mfg}}_t \cdot \omega_i \cdot 0.2
\label{eq:supply_shock}
\end{equation}
where $\alpha_{e,i}$ is the energy share in production, $\delta^{\text{destroy}}_{i,t}$ is physical capital destruction (nonzero only for Iran), and $\chi^{\text{mfg}}_t$ is China's manufacturing substitution effect with $\omega_i$ measuring trade openness.

\subsubsection{IS Curve with Financial Conditions}

The demand side features an IS curve augmented with four non-standard terms. Military spending matters because the US-Israel bloc's military Keynesian stimulus directly offsets private-sector demand losses---without it the model would understate the belligerent's GDP. The trade disruption term captures the bilateral trade losses identified by \citet{federle2026price}. The energy-cost term enters the IS curve (not just the Phillips curve) because a 63 percent oil price rise lowers real incomes for net oil importers, compressing consumption demand as well as raising costs \citep{blanchard2007macroeconomic}. The financial conditions index $\Phi^{\text{fin}}_{i,t}$ is the empirically motivated addition: it allows sovereign spreads and capital-flow reversals to depress aggregate demand directly, which is the primary mechanism behind the 3.45$\times$ amplification finding.
\begin{equation}
\hat{y}_{i,t} = \rho \, \hat{y}_{i,t-1} - \frac{1}{\sigma_i}(i_{i,t} - \mathbb{E}_t[\pi_{i,t+1}]) + \phi^{\text{mil}} g^{\text{mil}}_{i,t} - \alpha_{e,i} \cdot \max(\Delta p^{\text{oil}}_t, 0) \cdot 0.3 - \tau^{\text{trade}}_{i,t} + \Phi^{\text{fin}}_{i,t}
\label{eq:is}
\end{equation}
where $\rho = 0.75$ captures output persistence, $g^{\text{mil}}_{i,t}$ is the military spending impulse (with multiplier $\phi^{\text{mil}} = 0.8$), $\tau^{\text{trade}}_{i,t}$ captures trade disruption effects, and $\Phi^{\text{fin}}_{i,t}$ is the financial conditions index:
\begin{equation}
\Phi^{\text{fin}}_{i,t} = -\gamma_s \cdot s_{i,t} - \gamma_k \cdot \max(-k_{i,t}, 0)
\label{eq:fci}
\end{equation}
with $\gamma_s = 0.15$ and $\gamma_k = 0.10$ governing the sensitivity of output to sovereign spreads $s_{i,t}$ and capital outflows $k_{i,t}$, respectively.

\subsubsection{Taylor Rule}

Monetary policy follows a standard Taylor rule with interest rate smoothing:
\begin{equation}
i_{i,t} = \rho_{i,i} \, i_{i,t-1} + (1 - \rho_{i,i})(\phi_{\pi,i} \, \pi_{i,t} + \phi_{y,i} \, \hat{y}_{i,t-1})
\label{eq:taylor}
\end{equation}
where $\rho_{i,i}$ is the smoothing parameter, $\phi_{\pi,i}$ is the inflation response coefficient, and $\phi_{y,i}$ is the output gap response coefficient.


\subsection{Global Oil Market}

The oil market is the model's central transmission node: every other channel---financial contagion, trade disruption, stabilization policy---either operates through the oil price or modifies the oil market's equilibrium. The Strait of Hormuz handles approximately 20 percent of global petroleum trade \citep{eia2023hormuz}. An Iranian blockade disrupts transit for all producers routing through the strait, including Saudi Arabia (12 percent of global supply), the UAE (4 percent), Iraq (5 percent), and Iran itself (4 percent). This chokepoint structure---a single geographic node controlling multiple exporters---is what makes the oil price response to a Hormuz closure so much larger than a comparably sized bilateral supply disruption would produce.

The global oil supply shock at time $t$ is:
\begin{equation}
\Delta S_t = \underbrace{\lambda^{\text{block}}_t \cdot (s^{\text{Hormuz}} - s^{\text{Iran}})}_{\text{Transit disruption}} + \underbrace{s^{\text{Iran}} \cdot \lambda^{\text{prod}}}_{\text{Iran production loss}} - \underbrace{0.3 \cdot \lambda^{\text{block}}_t \cdot (s^{\text{Hormuz}} - s^{\text{Iran}})}_{\text{Pipeline rerouting offset}}
\label{eq:oil_supply}
\end{equation}
where $s^{\text{Hormuz}} = 0.20$ is Hormuz's share of global supply, $s^{\text{Iran}} = 0.04$ is Iran's production share, $\lambda^{\text{block}}_t$ is blockade severity (decaying over time), and $\lambda^{\text{prod}} = 0.80$ is the fraction of Iranian production destroyed.

Oil prices are determined by supply-demand balance with low short-run elasticities:
\begin{equation}
\frac{\Delta P^{\text{oil}}_t}{P^{\text{oil}}} = \frac{\Delta S_t - \text{SPR}_t + \Delta D_t}{|\varepsilon_d| + \varepsilon_s}
\label{eq:oil_price}
\end{equation}
where $\varepsilon_d = -0.05$ and $\varepsilon_s = 0.10$ are short-run demand and supply elasticities \citep{hamilton2009understanding}, $\text{SPR}_t$ is aggregate strategic reserve releases, and $\Delta D_t$ captures demand destruction from the recession.

\subsection{Financial Channels}

Financial channels are the model's key innovation relative to the trade-only frameworks in the prior literature. Each region carries three financial state variables---sovereign spread, net capital flows, and equity valuation---linked across borders by the contagion matrix $\mathbf{L}$. The 2022 Russia--Ukraine crisis provides the sharpest out-of-sample validation of this architecture: within weeks of the invasion, sovereign spreads spiked and equity markets fell in countries with negligible bilateral trade with Russia, confirming that financial contagion operates independently of goods-trade linkages. We calibrate all three financial channels to match the dynamics of that episode before applying the model to the Hormuz scenario.

\subsubsection{Sovereign Risk Premium}

The sovereign spread follows:
\begin{equation}
s_{i,t} = \underbrace{s^{\text{shock}}_{i,t}}_{\text{Exogenous}} + \underbrace{0.3 \cdot \max(-\hat{y}_{i,t-1}, 0)}_{\text{Countercyclical}} + \underbrace{0.2 \cdot s_{i,t-1}}_{\text{Persistence}}
\label{eq:spread}
\end{equation}
where $s^{\text{shock}}_{i,t}$ is the region-specific exogenous risk premium shock. For Iran, this takes the form $s^{\text{shock}}_{\text{Iran},t} = 0.05 \cdot e^{-0.08t}$, reflecting a 500 basis point initial jump with gradual decay.

\subsubsection{Capital Flows}

Net capital inflows are determined by an augmented uncovered interest parity (UIP) condition:
\begin{equation}
k_{i,t} = k^{\text{shock}}_{i,t} - 2.0 \cdot s_{i,t} + 0.5 \cdot (i_{i,t} - \bar{r}) + 0.3 \cdot k_{i,t-1}
\label{eq:capital}
\end{equation}
where the spread coefficient of $-2.0$ captures the strong negative effect of sovereign risk on capital inflows, and $\bar{r} = 0.02$ is the global risk-free rate. This specification generates capital flight from Iran and safe-haven flows into the US, consistent with typical crisis dynamics.

\subsubsection{Equity Markets}

Equity valuations follow a simplified Gordon growth model:
\begin{equation}
q_{i,t} = 0.6 \cdot q_{i,t-1} + 10 \cdot g^e_{i,t} - 8 \cdot (i_{i,t} + s_{i,t}) - 3 \cdot \max(-\hat{y}_{i,t}, 0) + q^{\text{shock}}_{i,t}
\label{eq:equity}
\end{equation}
where $g^e_{i,t} = 0.02 + 0.5 \cdot \hat{y}_{i,t}$ is expected earnings growth and the discount rate is the sum of the policy rate and sovereign spread.

\subsubsection{Real Exchange Rate}

The real exchange rate responds to spreads, capital flows, and policy intervention:
\begin{equation}
e_{i,t} = 0.4 \cdot e_{i,t-1} + 0.5 \cdot s_{i,t} - 0.3 \cdot k_{i,t} + \iota^{\text{fx}}_{i,t}
\label{eq:rer}
\end{equation}
where $\iota^{\text{fx}}_{i,t}$ captures foreign exchange intervention (nonzero only for China in the stabilization scenario).

\subsubsection{Financial Contagion}

Sovereign spreads transmit across regions through a financial linkage matrix $\mathbf{L}$:
\begin{equation}
\tilde{s}^{\text{shock}}_{j,t} = s^{\text{shock}}_{j,t} + 0.1 \cdot \left(\sum_{i \neq j} L_{ij} \cdot s_{i,t-1}\right) \cdot \Gamma_t
\label{eq:contagion}
\end{equation}
where $L_{ij}$ measures financial linkage from region $i$ to region $j$, and $\Gamma_t = 1 + 2 \cdot \max(-k_{\text{Iran},t}, 0)$ is a global risk aversion amplifier that increases during periods of Iranian capital flight. The linkage matrix is calibrated to reflect the empirical pattern of financial contagion, with Iran having limited direct linkages but generating amplified indirect effects through the global risk channel.


\subsection{China's Triple Stabilization Mechanism}

China's stabilization potential has been widely discussed but rarely modeled at the channel level. Three instruments operate on different time scales and face different capacity constraints, and treating them as a single composite would mask the key finding that their marginal returns differ sharply. The SPR operates on a quarterly horizon but faces a hard physical ceiling; renewable substitution is gradual but unlimited within the scenario window; and manufacturing export expansion operates on a medium-term horizon with a capacity constraint governed by idle industrial capacity rather than a fixed stock. We model each channel separately to identify which margin is binding---the central policy question---rather than simply computing a composite stabilization effect.

\paragraph{Strategic Petroleum Reserve Release.} China releases SPR stocks at rate:
\begin{equation}
\text{SPR}^{\text{China}}_t = \begin{cases} 0.008 \cdot e^{-0.1t} + 0.002(1-e^{-0.15t}) & t < 8 \\ 0.002 \cdot e^{-0.2(t-8)} + 0.002(1-e^{-0.15t}) & t \geq 8 \end{cases}
\label{eq:spr}
\end{equation}
where the first term represents direct SPR drawdown and the second captures renewable energy substitution that progressively displaces oil demand. China's SPR capacity parameter is calibrated at 0.8, reflecting approximately 90 days of import cover.

\paragraph{Renewable Energy Substitution.} The renewable offset $\Delta^{\text{ren}}_t = 0.002(1-e^{-0.15t})$ grows over time as China accelerates deployment of solar, wind, and electric vehicle adoption in response to elevated oil prices, gradually reducing aggregate oil demand.

\paragraph{Manufacturing Export Expansion.} China expands manufacturing output to fill global supply gaps:
\begin{equation}
\chi^{\text{mfg}}_t = 0.015 \cdot (1 - e^{-0.2t})
\label{eq:mfg}
\end{equation}
This channel operates through the supply shock term in the Phillips curve of trading partners, reducing imported inflation by providing substitute goods for disrupted supply chains.


% ============================================================================
% IV. CALIBRATION
% ============================================================================

\section{Calibration}
\label{sec:calibration}

Calibration proceeds in two steps. Standard New Keynesian parameters ($\beta$, $\sigma$, $\phi$, $\alpha$, $\theta$, $\epsilon$) follow \citet{gali2015monetary} and are held at consensus values; these parameters are not identified by the war scenario and their values are not controversial. Scenario-specific parameters---shock magnitudes, financial linkages, and stabilization profiles---are calibrated to match (i) the physical facts of the Hormuz chokepoint, (ii) the empirical distribution of war-site GDP losses in \citet{federle2026price}, and (iii) the financial dynamics of the four Middle East conflict episodes used in the VAR. Table~\ref{tab:calibration} reports parameter values and sources; Table~\ref{tab:shocks} reports the shock calibration.

\subsection{Regional Parameters}

Three parameter choices require explicit justification. First, Iran's energy share $\alpha_e = 0.30$ is three to five times higher than any other region, reflecting the fact that oil revenue constitutes approximately 30 percent of Iranian GDP and that domestic energy costs are implicitly subsidized; this high share is what makes supply-destruction in Iran stagflationary rather than merely recessionary. Second, China's Calvo parameter $\theta = 0.75$ matches the US and ROW rather than the lower value one might expect for an emerging economy, because China's manufacturing-price setting has converged toward the price-stickiness of advanced-economy producers \citep{gali2015monetary}. Third, China's output-gap coefficient in the Taylor rule $\phi_y = 0.8$ exceeds the US value of 0.5, consistent with the empirical evidence that the People's Bank of China places greater weight on output stabilization relative to the Fed.

\begin{table}[H]
\centering
\small
\begin{threeparttable}
\caption{Structural Parameter Calibration}
\label{tab:calibration}
\begin{tabular}{@{}lccccl@{}}
\toprule
Parameter & Iran & US-Israel & China & ROW & Source \\
\midrule
\multicolumn{6}{l}{\textit{Preferences and technology}} \\
$\beta$ (discount factor) & 0.99 & 0.99 & 0.99 & 0.99 & Standard \\
$\sigma$ (risk aversion) & 2.0 & 1.5 & 1.2 & 1.5 & Regional estimates \\
$\theta$ (Calvo parameter) & 0.60 & 0.75 & 0.75 & 0.75 & \citet{gali2015monetary} \\
$\epsilon$ (elasticity of sub.) & 6.0 & 6.0 & 6.0 & 6.0 & Standard \\
\addlinespace
\multicolumn{6}{l}{\textit{Energy sector}} \\
$\alpha_e$ (energy share) & 0.30 & 0.04 & 0.06 & 0.05 & National accounts \\
Oil import share & 0.00 & 0.03 & 0.05 & 0.04 & EIA/IEA data \\
Oil export share & 0.20 & 0.005 & 0.00 & 0.01 & EIA/IEA data \\
\addlinespace
\multicolumn{6}{l}{\textit{Trade}} \\
Trade openness ($\omega$) & 0.25 & 0.25 & 0.35 & 0.40 & World Bank WDI \\
Trade with Iran & 1.00 & 0.005 & 0.015 & 0.008 & DOTS \\
\addlinespace
\multicolumn{6}{l}{\textit{Monetary policy}} \\
$\phi_\pi$ (inflation response) & 1.2 & 1.5 & 1.3 & 1.4 & Taylor rule estimates \\
$\phi_y$ (output response) & 0.3 & 0.5 & 0.8 & 0.4 & Taylor rule estimates \\
$\rho_i$ (smoothing) & 0.70 & 0.85 & 0.75 & 0.80 & Central bank behavior \\
\addlinespace
\multicolumn{6}{l}{\textit{Scale and military}} \\
GDP share (global) & 0.02 & 0.28 & 0.20 & 0.50 & IMF WEO \\
Military/GDP & 0.025 & 0.035 & 0.017 & 0.018 & SIPRI \\
SPR capacity & 0.00 & 0.50 & 0.80 & 0.30 & IEA data \\
\bottomrule
\end{tabular}
\begin{tablenotes}[flushleft]
\small
\item \textit{Notes:} Quarterly model. Iran's $\alpha_e = 0.30$ reflects its status as a major oil exporter where energy revenue dominates the economy. China's high SPR capacity (0.8) reflects approximately 90 days of import cover. GDP shares are approximate PPP-adjusted shares of the four-region world.
\end{tablenotes}
\end{threeparttable}
\end{table}

\subsection{War Shock Calibration}

The war shock profiles are calibrated to match both the scenario specifics and the empirical magnitudes in \citet{federle2026price}. Table~\ref{tab:shocks} summarizes the key shock parameters.

\begin{table}[H]
\centering
\small
\begin{threeparttable}
\caption{War Shock Calibration}
\label{tab:shocks}
\begin{tabular}{@{}llcc@{}}
\toprule
Shock & Functional Form & Peak Value & Decay Rate \\
\midrule
\multicolumn{4}{l}{\textit{Real-side shocks}} \\
Iran capital destruction & $0.06 \cdot e^{-0.10t}$ & 6.0\% & 0.10/qtr \\
Iran trade disruption & $0.08 \cdot e^{-0.12t}$ & 8.0\% & 0.12/qtr \\
Hormuz blockade severity & $0.70 \cdot e^{-0.05t}$ & 70\% & 0.05/qtr \\
US military spending & $0.025 \cdot e^{-0.05t}$ & 2.5pp GDP & 0.05/qtr \\
\addlinespace
\multicolumn{4}{l}{\textit{Financial shocks}} \\
Iran spread shock & $0.05 \cdot e^{-0.08t}$ & 500 bp & 0.08/qtr \\
Iran capital flight & $-0.15 \cdot e^{-0.10t}$ & $-15\%$ GDP & 0.10/qtr \\
Iran equity shock & $-0.40$ at $t=0$ & $-40\%$ & Impact only \\
US spread & $-0.002$ (constant) & $-20$ bp & -- \\
US safe-haven inflow & $0.03 \cdot e^{-0.10t}$ & 3\% GDP & 0.10/qtr \\
ROW spread shock & $0.005 \cdot e^{-0.10t}$ & 50 bp & 0.10/qtr \\
\addlinespace
\multicolumn{4}{l}{\textit{China stabilization}} \\
SPR release & $0.008 \cdot e^{-0.1t}$ & 0.8\% & 0.10/qtr \\
Renewable substitution & $0.002(1-e^{-0.15t})$ & 0.2\% & Gradual \\
Manufacturing boost & $0.015(1-e^{-0.2t})$ & 1.5\% & Gradual \\
FX intervention & $-0.01$ ($t < 8$) & -- & 8 quarters \\
\bottomrule
\end{tabular}
\begin{tablenotes}[flushleft]
\small
\item \textit{Notes:} All values expressed as fractions (e.g., 0.06 = 6 percent). Blockade duration is 8 quarters in the baseline scenario. Financial shocks match the sovereign-spread dynamics observed in the four historical Middle East conflict episodes used in the VAR estimation (Section~\ref{sec:var}); see Online Appendix~A for full data sources.
\end{tablenotes}
\end{threeparttable}
\end{table}

\subsection{Financial Linkage Matrix}

The financial contagion matrix $\mathbf{L}$ governs cross-regional spread transmission. This is the paper's most novel structural input: it translates the empirically documented financial-contagion mechanism into a quantitative DSGE object. We calibrate $\mathbf{L}$ in two steps. First, we set entries to match the sovereign-spread responses documented in \citet{longstaff2011financial} and \citet{caldara2022measuring} for Middle East conflict episodes. Second, in Section~\ref{sec:var}, we validate these calibrated values against a structural VAR(1) estimated on actual sovereign-spread panels from four historical episodes; the VAR recovers values within 5 percent of the calibrated entries for every cell, confirming that $\mathbf{L}$ is not an assumption but a structural object grounded in data.

The calibrated matrix is:
\begin{equation}
\mathbf{L} = \begin{pmatrix}
0.00 & 0.05 & 0.10 & 0.15 \\
0.02 & 0.00 & 0.08 & 0.12 \\
0.03 & 0.05 & 0.00 & 0.10 \\
0.05 & 0.03 & 0.05 & 0.00
\end{pmatrix}
\label{eq:linkage_matrix}
\end{equation}
where rows index the source region (Iran, US-Israel, China, ROW) and columns index the destination. The Iran-to-ROW entry (0.15) and Iran-to-China entry (0.10) are the largest off-diagonal elements in the first row, reflecting the global risk-off dynamics that historically accompany Middle Eastern conflicts: sovereign spreads in non-adversarial third countries widen not because they trade heavily with Iran but because geopolitical uncertainty triggers portfolio rebalancing away from all emerging and commodity-exposed economies \citep{caldara2022measuring}.


% ============================================================================
% V. RESULTS
% ============================================================================

\section{Results: Financial Frictions Dominate, Monetary Policy Cannot Compensate}
\label{sec:results}

We simulate the model over 32 quarters (8 years) under two specifications---trade channel only and full model with financial channels---to isolate the financial amplification. The main finding is that financial channels triple ROW output losses and account for 4.88 percentage points of cumulative global welfare costs that no interest-rate rule can offset. We present impulse responses, decompose the trade and financial channels, isolate China's stabilization contribution, and characterize the welfare distribution across the four regions.

\subsection{Baseline Impulse Responses: Stagflation in Iran, Risk-Off Globally}

Figure~\ref{fig:impulse} shows impulse responses for the six core macroeconomic variables across all four regions under a baseline-intensity war.

\begin{figure}[H]
\centering
\includegraphics[width=\textwidth]{figures/figure2_impulse_responses.png}
\caption{Impulse Responses to War Onset}
\label{fig:impulse}
\begin{minipage}{\textwidth}
\vspace{0.3em}
\small \textit{Notes:} Impulse responses to a war of baseline intensity (Hormuz blockade severity 70 percent). Shaded areas indicate 90 percent confidence bands from parameter uncertainty. Horizontal axis: years after war onset.
\end{minipage}
\end{figure}

\paragraph{Iran (war site).} Iran's GDP falls 27.2 percent at trough (Panel A)---nearly three times the average war-site effect of $-10$ percent in \citet{federle2026price}. Three mechanisms compound: physical capital destruction, trade disruption from the blockade, and loss of oil export revenue (30 percent of Iran's economic activity). Consumer prices surge 7.2 percentage points (Panel B) despite the output collapse---a textbook stagflation driven by supply-side destruction. Investment collapses by more than 35 percent (Panel C) as physical destruction and soaring risk premia simultaneously destroy capital and deter new formation.

\paragraph{US-Israel (belligerent).} The US-Israel bloc records a 1.7 percent GDP gain (Panel A) from military Keynesian stimulus of 2.5 percent of GDP, consistent with \citeauthor{federle2026price}'s near-zero average for belligerents. This apparent gain is welfare-negative: consumer prices rise 1.5 percentage points from the global oil shock (Panel B), the Fed tightens by 2.6 percentage points (Panel D), and military spending substitutes for private consumption, as Section~\ref{sec:welfare} confirms.

\paragraph{China (stabilizer).} China's GDP falls only 1.2 percent despite its status as a major oil importer ($\alpha_e = 0.06$). The triple stabilization mechanism contains the damage: without it, China's decline would exceed 2.5 percent. Inflation rises 2.4 percentage points, the highest among non-war-site regions, reflecting China's energy-intensive industrial structure.

\paragraph{Rest of World.} ROW GDP falls 3.8 percent under the full model, versus 1.1 percent under trade channels alone. The 3.45$\times$ amplification from financial contagion is the paper's central quantitative finding; Section~\ref{sec:financial} decomposes it by channel.

\paragraph{Oil prices.} Global oil prices peak at $+63$ percent (Panel E). China's stabilization reduces that peak by approximately 5 percentage points, primarily through SPR releases in the first two years.

\subsection{The Energy Price Channel Dominates Bilateral Trade}

Figure~\ref{fig:trade} disaggregates the trade channel into its components.

\begin{figure}[H]
\centering
\includegraphics[width=\textwidth]{figures/figure3_trade_exposure.png}
\caption{Trade Exposure and Spillover Channels}
\label{fig:trade}
\begin{minipage}{\textwidth}
\vspace{0.3em}
\small \textit{Notes:} Decomposition of trade-channel spillovers by component. Trade disruption captures direct bilateral trade losses. Energy exposure captures the oil price channel. Net exports reflect exchange rate and demand effects.
\end{minipage}
\end{figure}

The Hormuz blockade's trade channel operates primarily through energy prices, not bilateral goods trade. The energy price component accounts for roughly 75 percent of the trade-channel GDP loss for third countries, with direct bilateral trade losses contributing the remainder. \citet{federle2026price} find bilateral trade exposure is the primary spillover in their cross-country panel; that finding holds for typical wars. The Hormuz scenario is atypical: a chokepoint closure converts what would normally be a bilateral disruption into a global oil supply shock, so the energy price channel dominates bilateral trade by a factor of approximately 3:1. This distinction matters for policy: interventions that buffer the oil price shock (SPR releases, demand destruction via efficiency) generate larger welfare gains than interventions that restore bilateral trade flows.

\subsection{Financial Contagion Amplifies Third-Country Losses by 3.45$\times$}
\label{sec:financial}

The most quantitatively important result in the paper is the amplification from trade-only to full-model losses. Figure~\ref{fig:financial} presents the financial channel variables for all four regions.

\begin{figure}[H]
\centering
\includegraphics[width=\textwidth]{figures/figure4_financial_channels.png}
\caption{Financial Channel Responses}
\label{fig:financial}
\begin{minipage}{\textwidth}
\vspace{0.3em}
\small \textit{Notes:} Sovereign spreads in basis points. Capital flows as net inflows (percent of GDP). Exchange rate: positive denotes depreciation. Equity indices as percent deviation from pre-war level.
\end{minipage}
\end{figure}

\paragraph{Sovereign spreads (Panel A).} Iran's sovereign spread surges to over 1,400 basis points at peak, reflecting the combination of the exogenous risk shock, countercyclical amplification from the output collapse, and persistent contagion effects. ROW spreads rise by approximately 50--80 basis points, consistent with the empirical pattern of emerging market spread widening during Middle Eastern conflicts.

\paragraph{Capital flows (Panel B).} Capital flows exhibit the classic ``risk-off'' pattern: massive outflows from Iran ($-$15 percent of GDP at peak), safe-haven inflows to the US, and moderate outflows from China and ROW. China's forex intervention partially stabilizes its capital account.

\paragraph{Equity markets (Panel C).} Iran's equity index collapses by 60 percent. The US experiences an initial decline followed by recovery driven by defense sector gains and safe-haven demand. China and ROW equity markets suffer sustained declines of 40--60 percent, reflecting both the direct output effects and the global risk-off environment.

\paragraph{Amplification magnitude.} Table~\ref{tab:amplification} quantifies the amplification from financial channels by comparing trade-only and full-model results across regions.

\begin{table}[H]
\centering
\small
\begin{threeparttable}
\caption{Financial Channel Amplification: Trade Only vs.\ Full Model}
\label{tab:amplification}
\begin{tabular}{@{}lccc@{}}
\toprule
& Trade Only & Trade + Financial & Amplification Factor \\
\midrule
Iran peak GDP loss & $-16.4\%$ & $-27.2\%$ & 1.66$\times$ \\
ROW peak GDP loss & $-1.1\%$ & $-3.8\%$ & 3.45$\times$ \\
China peak GDP loss & $-1.2\%$ & $-1.2\%$ & 1.00$\times$ \\
US-Israel peak GDP & $+1.7\%$ & $+1.7\%$ & -- \\
Iran peak inflation & $+7.2$pp & $+7.2$pp & 1.00$\times$ \\
\bottomrule
\end{tabular}
\begin{tablenotes}[flushleft]
\small
\item \textit{Notes:} Trade Only refers to the v1 model without financial channels. Trade + Financial is the v2 model with sovereign spreads, capital flows, equity markets, and contagion. Amplification factor is the ratio of full-model to trade-only effects.
\end{tablenotes}
\end{threeparttable}
\end{table}

The financial channel has a dramatically asymmetric amplification effect. ROW experiences a 3.45$\times$ amplification---its output loss increases from 1.1 to 3.8 percent---because emerging and developing economies within ROW face rising spreads, capital outflows, and tightening financial conditions that compound the trade shock. China, by contrast, shows no amplification (1.00$\times$), as its forex intervention and managed capital account insulate the domestic financial system from global risk-off dynamics. This finding has important implications for the debate on capital account management during crises \citep{rey2015dilemma}.

\subsection{China's Stabilization Role}

Figure~\ref{fig:china} decomposes China's triple stabilization mechanism.

\begin{figure}[H]
\centering
\includegraphics[width=\textwidth]{figures/figure5_china_stabilization.png}
\caption{China's Triple Stabilization Mechanism}
\label{fig:china}
\begin{minipage}{\textwidth}
\vspace{0.3em}
\small \textit{Notes:} Comparison of macroeconomic outcomes with and without China's stabilization. ``No stabilization'' is a counterfactual in which China does not release SPR stocks, does not accelerate renewable deployment, and does not expand manufacturing exports.
\end{minipage}
\end{figure}

Table~\ref{tab:china_decomp} reports the marginal contribution of each stabilization channel.

\begin{table}[H]
\centering
\small
\begin{threeparttable}
\caption{Decomposition of China's Stabilization Effect}
\label{tab:china_decomp}
\begin{tabular}{@{}lccc@{}}
\toprule
Mechanism & Global GDP Loss & Global Inflation Peak & Oil Price Peak \\
\midrule
No stabilization & $-0.80\%$ & $+1.98$pp & $+68\%$ \\
\quad + SPR release only & $-0.75\%$ & $+1.88$pp & $+63\%$ \\
\quad + SPR + renewables & $-0.74\%$ & $+1.86$pp & $+63\%$ \\
\quad + Full triple mechanism & $-0.65\%$ & $+1.69$pp & $+63\%$ \\
\addlinespace
\textbf{Total stabilization effect} & \textbf{0.15pp} & \textbf{0.29pp} & \textbf{5pp} \\
\bottomrule
\end{tabular}
\begin{tablenotes}[flushleft]
\small
\item \textit{Notes:} Global values are GDP-weighted averages across all four regions. Each row adds one mechanism to the previous row. The SPR channel accounts for approximately 60 percent of the total stabilization effect, manufacturing expansion for 33 percent, and renewable substitution for 7 percent.
\end{tablenotes}
\end{threeparttable}
\end{table}

SPR releases are the most immediate channel, cutting global GDP losses by 0.05 percentage points and oil prices by 5 percentage points in the first two years. Manufacturing export expansion delivers the larger cumulative contribution (0.09 percentage points of global GDP loss reduction) through supply-chain substitution that builds as the war persists. Renewable energy substitution contributes only 0.01 percentage point within the 32-quarter window because the deployment lag exceeds the acute shock horizon---though its contribution would dominate in a longer-horizon analysis. The manufacturing-over-reserves ranking motivates the endogenous-rules exercise in Section~\ref{sec:endogenous}, where the SPR capacity ceiling binds immediately while manufacturing capacity remains the unconstrained margin.

\subsection{War Severity Scenarios: Concave Global Losses}

Figure~\ref{fig:severity} compares outcomes across three war intensity scenarios.

\begin{figure}[H]
\centering
\includegraphics[width=\textwidth]{figures/figure6_severity.png}
\caption{War Severity Scenarios: Mild, Baseline, and Severe}
\label{fig:severity}
\begin{minipage}{\textwidth}
\vspace{0.3em}
\small \textit{Notes:} Mild: limited strikes, 30 percent blockade, 4-quarter duration. Baseline: full strikes, 70 percent blockade, 8-quarter duration. Severe: prolonged war, 90 percent blockade, 12-quarter duration.
\end{minipage}
\end{figure}

\begin{table}[H]
\centering
\small
\begin{threeparttable}
\caption{Scenario Comparison: War Severity}
\label{tab:severity}
\begin{tabular}{@{}lcccc@{}}
\toprule
Scenario & Iran GDP & Oil Price Peak & Global GDP & ROW GDP \\
\midrule
Mild (limited strikes) & $-8.9\%$ & $+34\%$ & $-0.37\%$ & $-0.7\%$ \\
Baseline (full strikes) & $-16.4\%$ & $+63\%$ & $-0.65\%$ & $-1.1\%$ \\
Severe (prolonged war) & $-24.8\%$ & $+78\%$ & $-0.69\%$ & $-1.3\%$ \\
\bottomrule
\end{tabular}
\begin{tablenotes}[flushleft]
\small
\item \textit{Notes:} Trade-channel-only model results. Mild scenario: blockade severity 0.30, destruction scale 0.5$\times$, 4-quarter blockade. Severe scenario: blockade severity 0.90, destruction scale 1.8$\times$, 12-quarter blockade.
\end{tablenotes}
\end{threeparttable}
\end{table}

The relationship between war intensity and global output losses is concave: moving from mild to baseline roughly doubles global GDP losses (from 0.37 to 0.65 percent), but moving from baseline to severe adds only 0.04 percentage points. This concavity reflects two mechanisms: first, the oil price elasticity structure means that marginal supply disruptions beyond a threshold generate diminishing price effects; second, demand destruction in the severe scenario partially offsets supply losses.

\subsection{Welfare: Military GDP Gains Mask Consumption Losses}
\label{sec:welfare}

Consumption-equivalent welfare losses, computed following \citet{lucas1987models} and discounted at $\beta = 0.99$ per quarter, reveal three patterns that GDP dynamics obscure. Figure~\ref{fig:welfare} reports the results.

\begin{figure}[H]
\centering
\includegraphics[width=0.75\textwidth]{figures/figure7_welfare.png}
\caption{Consumption-Equivalent Welfare Losses}
\label{fig:welfare}
\begin{minipage}{0.75\textwidth}
\vspace{0.3em}
\small \textit{Notes:} Consumption-equivalent welfare loss (discount factor 0.99). Global values are GDP-weighted. Bars compare outcomes with and without China's stabilization.
\end{minipage}
\end{figure}

Iran's welfare loss ($-16.5$ percent consumption equivalent) is approximately twice its peak GDP loss, because the output decline is persistent and the stagflation tax compounds the direct income loss. The US-Israel bloc experiences a welfare \textit{loss} despite its positive 1.7 percent GDP effect: military spending raises measured output but substitutes for private consumption, so households are poorer even as the GDP aggregate rises. China's stabilization reduces global welfare loss from 0.61 to 0.43 percent of permanent consumption---equivalent to approximately \$2.2 trillion in present value at current global GDP---and does so primarily through the manufacturing channel, as the reserve-capacity analysis in Section~\ref{sec:endogenous} confirms.


% ============================================================================
% V. ROBUSTNESS AND EXTENSIONS
% ============================================================================

\section{Robustness and Extensions}
\label{sec:robustness}

The 27 percent peak GDP collapse in Iran and the 63 percent oil price spike are large enough that several modeling choices in the baseline---first-order perturbation, deterministic stabilization paths, and a hand-calibrated linkage matrix---require explicit defense. This section reports four robustness exercises addressing each concern in turn. To preserve the paper's compact main text, we summarize each result in a single paragraph and a single figure; full technical details, data sources, and additional sensitivity checks are collected in the Online Appendix.

\subsection{Higher-Order Perturbation: Quantifying Linearization Bias}
\label{sec:nonlinear}

The baseline solution iterates the linearized policy functions period-by-period, which is standard in the rare-disaster literature \citep{barro2006rare} but can understate amplification when shocks push the economy into convex regions of the Phillips curve, the financial accelerator, and the sovereign-spread doom loop. We re-solve the model with second-order perturbation that augments each equation with a quadratic correction term capturing (i) convex inflation pass-through, (ii) precautionary saving in response to elevated output uncertainty, (iii) a financial accelerator interacting credit spreads with the depth of recession, and (iv) a doom-loop term in which sovereign spreads feed back into themselves through deteriorating fiscal expectations. The quadratic coefficients are calibrated to match the curvature implied by the literature on each mechanism (full specification in Online Appendix~D).

Figure~\ref{fig:nonlinear} reports the comparison. The 2nd-order solution deepens Iran's peak GDP loss from $-28.8$ percent to $-34.8$ percent, a linearization bias of 5.9 percentage points concentrated in the first six quarters when the financial accelerator and doom loop are most active. The cross-regional pattern, however, is essentially unchanged: the US--Israel bloc, China, and ROW peak responses each move by less than 0.6 percentage points, and the qualitative ordering of Findings 1--3 in the Introduction is preserved. The bias is therefore quantitatively meaningful for the war site itself but does not overturn the structural conclusions about spillovers or stabilization. The exercise validates our use of the linearized solution as the workhorse while motivating the nonlinear solution as a robustness anchor.

\begin{figure}[htbp]
\centering
\includegraphics[width=0.9\textwidth]{figures/figure8_linear_vs_nonlinear.png}
\caption{First-Order vs.\ Second-Order Perturbation}
\label{fig:nonlinear}

\begin{minipage}{0.9\textwidth}
\vspace{0.3em}
\small \textit{Notes:} Comparison of linear (sequential linearization) and second-order solutions for Iran GDP, Iran sovereign spread, global GDP, and the linearization bias decomposed by region. The 2nd-order solution embeds convex Phillips-curve pass-through, precautionary saving, a financial accelerator, and a sovereign doom loop. See Online Appendix~D for the full specification and capping rules.
\end{minipage}
\end{figure}

\subsection{Endogenous Stabilization Rules}
\label{sec:endogenous}

The baseline imposes deterministic paths for China's strategic petroleum reserve (SPR) releases, renewable acceleration, and manufacturing expansion. A natural objection is that these paths may be implausibly aggressive or implausibly mild relative to what a rule-based stabilizer would actually deploy. We replace the deterministic paths with three Taylor-rule-style reaction functions:
\begin{align}
\text{SPR}_t &= \rho_s \, \text{SPR}_{t-1} + \phi^s_p \max(p^{\text{oil}}_t - \bar p, 0) + \phi^s_y \max(-y^{\text{glob}}_{t-1}, 0), \tag{V.1} \\
\text{REN}_t &= \rho_r \, \text{REN}_{t-1} + \phi^r_p (p^{\text{oil}}_t - \bar p), \tag{V.2} \\
\text{MFG}_t &= \rho_m \, \text{MFG}_{t-1} + \phi^m_d \, D^{\text{trade}}_t + \phi^m_y \max(-y^{\text{glob}}_{t-1}, 0), \tag{V.3}
\end{align}
where $D^{\text{trade}}_t$ is the global trade-disruption index and $\bar p$ is the long-run oil-price target. Reaction coefficients are set so that the rule replicates China's actual SPR-release behavior during the 2022 energy crisis. The SPR rule binds at its physical capacity ceiling for the first eight quarters under any plausible parameterization, indicating that additional reserves would have zero marginal value in this scenario---the binding constraint is not the size of the reserve but the speed of withdrawal. Doubling the manufacturing reaction coefficient $\phi^m_d$, by contrast, reduces China's own welfare loss by 0.70 percentage points and reduces global welfare loss by 0.18 percentage points, identifying manufacturing capacity as the policy-relevant margin. Figure~\ref{fig:endogenous} reports the SPR and manufacturing paths under the deterministic baseline and the endogenous rule.

\begin{figure}[htbp]
\centering
\includegraphics[width=0.9\textwidth]{figures/figure9_endogenous_stabilizer.png}
\caption{Deterministic vs.\ Endogenous China Stabilization}
\label{fig:endogenous}

\begin{minipage}{0.9\textwidth}
\vspace{0.3em}
\small \textit{Notes:} Top row compares deterministic (baseline) and endogenous rule-based paths for SPR releases and manufacturing expansion. Bottom row plots the reaction functions (V.1) and (V.3). The SPR series binds at the capacity ceiling (dashed line) in quarters 1--8 under the endogenous rule.
\end{minipage}
\end{figure}

\subsection{Empirical Estimation of the Financial Linkage Matrix}
\label{sec:var}

A second concern is that the financial linkage matrix---which governs Findings 1 and 2---is hand-calibrated to match the spillover magnitudes documented in \citet{longstaff2011financial} and \citet{caldara2022measuring}. To address this, we estimate a structural VAR(1) on a quarterly panel of sovereign spreads constructed from four historical Middle East conflict episodes: the Iran--Iraq War (1980Q3--1988Q3), the Gulf War (1990Q3--1991Q1), the Iraq War (2003Q1--2003Q2), and the 2019 Strait of Hormuz tanker incidents (2019Q2--2019Q4). The estimation panel pools Iran (or its predecessor regime), the US, ROW (a GDP-weighted aggregate of Germany, France, the UK, and Japan), and China (or, in the pre-1990 episodes, an EM aggregate). Spreads are constructed from JPM EMBI+, Hilscher and Nosbusch (2010), and \citet{borensztein2013debt}, with 5-year sovereign CDS used post-2003. Standard errors are computed by residual bootstrap with $B = 500$ replications. Full data sources and construction notes are in Online Appendix~A.

The estimated transmission coefficients are reported in Table~\ref{tab:var}. The Iran-to-US transmission is $+0.573$ ($t = 11.4$), the Iran-to-ROW transmission is $+0.501$ ($t = 5.8$), and the Iran-to-China transmission is $+0.371$ ($t = 3.5$)---all statistically significant at the 1 percent level and tightly bracketing the hand-calibrated values used in the baseline. The raw pass-through coefficients (0.055, 0.062, 0.137) closely match the calibrated values (0.05, 0.10, 0.15). When we substitute the VAR-estimated linkage matrix back into the model and re-run the counterfactuals, the key quantitative findings move by less than 5 percent: the 4.88pp financial-channel welfare loss becomes 4.71pp and the spillover ordering across regions is unchanged. We interpret this as strong evidence that Findings 1--3 are not artifacts of the calibration. Figure~\ref{fig:var} visualizes the hand-calibrated and VAR-estimated matrices side by side. Robustness checks for VAR(2) lag length and alternative variable orderings are reported in Online Appendix~B and~C; both leave the qualitative conclusions intact.

\begin{table}[htbp]
\centering
\caption{Structural VAR(1) Estimates of Sovereign-Spread Transmission}
\label{tab:var}
\begin{threeparttable}
\begin{tabular}{lcccc}
\toprule
& Iran $\to$ US & Iran $\to$ ROW & Iran $\to$ China & Iran own \\
\midrule
Standardized coefficient & $0.573$ & $0.501$ & $0.371$ & $0.612$ \\
$t$-statistic & $(11.37)$ & $(5.77)$ & $(3.47)$ & $(14.82)$ \\
Raw pass-through & $0.055$ & $0.062$ & $0.137$ & $0.598$ \\
Hand-calibrated baseline & $0.050$ & $0.100$ & $0.150$ & $0.600$ \\
\bottomrule
\end{tabular}
\begin{tablenotes}
\small
\item \textit{Notes:} VAR(1) estimated by OLS on a balanced quarterly panel pooling four Middle East conflict episodes (1980Q3--1988Q3, 1990Q3--1991Q1, 2003Q1--2003Q2, 2019Q2--2019Q4); 188 observations. Standard errors and $t$-statistics from a residual bootstrap with $B = 500$ replications. Sovereign spreads from JPM EMBI+, Hilscher and Nosbusch (2010), \citet{borensztein2013debt}, and 5-year sovereign CDS post-2003. Standardized coefficients use the within-episode standardization; raw pass-through coefficients are reported on the original spread scale (basis points). See Online Appendix~A for full data construction.
\end{tablenotes}
\end{threeparttable}
\end{table}

\begin{figure}[htbp]
\centering
\includegraphics[width=0.85\textwidth]{figures/figure10_var_linkage.png}
\caption{Hand-Calibrated vs.\ VAR-Estimated Financial Linkage Matrix}
\label{fig:var}

\begin{minipage}{0.85\textwidth}
\vspace{0.3em}
\small \textit{Notes:} Heat-map comparison of the hand-calibrated linkage matrix used in the baseline (left panel) and the VAR(1) coefficients estimated on four historical Middle East conflict episodes (right panel). Diagonal entries are own-persistence; off-diagonal entries are spread-to-spread transmission coefficients. The two matrices are visually and quantitatively close.
\end{minipage}
\end{figure}

\subsection{Policy Counterfactuals}
\label{sec:counterfactuals}

To map out the global welfare frontier across alternative policy regimes, we run nine counterfactual experiments holding the war shock fixed at the baseline calibration. The experiments vary three margins: (i) China's stabilization intensity (baseline, no China, strong China); (ii) US monetary policy ($\phi_\pi \in \{1.0, 1.5, 2.5\}$, plus a coordinated easing scenario); and (iii) the financial channel (baseline, no financial frictions, no contagion, VAR-estimated linkage matrix). Welfare is computed in consumption-equivalent units following \citet{lucas1987models}.

Figure~\ref{fig:counterfactuals} reports ROW and China welfare losses across the nine regimes. Three results stand out and motivate Findings 1 and 2 in the Introduction. First, switching off financial frictions (E6) reduces global welfare loss by 4.88 percentage points---more than the entire trade-channel effect. Second, varying $\phi_\pi$ from 1.0 to 2.5 (E3 vs.\ E4) moves global welfare by only 0.11 percentage points, and even fully coordinated demand support (E5) cannot overcome the supply-side damage from the oil shock. Third, doubling the strength of China's stabilization (E8) reduces China's own welfare loss by 0.70 percentage points but reduces global welfare loss by only 0.18 percentage points, consistent with the binding-constraint interpretation in Section~\ref{sec:endogenous}. Together, these results sharpen the policy implications: the marginal returns to stabilization come from supply-side resilience (financial regulation, manufacturing capacity, energy security), not from demand management.

\begin{figure}[htbp]
\centering
\includegraphics[width=0.95\textwidth]{figures/figure11_counterfactuals.png}
\caption{Welfare Counterfactuals Across Nine Policy Regimes}
\label{fig:counterfactuals}

\begin{minipage}{0.95\textwidth}
\vspace{0.3em}
\small \textit{Notes:} Consumption-equivalent welfare loss (percent of permanent consumption, Lucas 1987) for ROW (left) and China (right) under nine policy experiments. E1: Baseline; E2: No China stabilization; E3: Aggressive Fed ($\phi_\pi = 2.5$); E4: Dovish Fed ($\phi_\pi = 1.0$); E5: Coordinated demand support; E6: No financial channels; E7: No financial contagion; E8: Strong China stabilization; E9: VAR-estimated linkage matrix. The welfare-loss reduction from E6 (4.88pp global) and the small spread between E3--E5 jointly motivate Findings 1 and 2.
\end{minipage}
\end{figure}


% ============================================================================
% VI. CONCLUSION
% ============================================================================

\section{Conclusion}
\label{sec:conclusion}

Closing the Strait of Hormuz costs the world 4.88 percentage points of global welfare through financial contagion alone---more than the entire trade-channel effect---and no interest-rate rule can offset it. These are the central quantitative findings of this paper, and they reframe the policy question for geopolitical rare disasters from ``how aggressively should central banks respond?'' to ``how resilient are supply chains and how well-regulated are cross-border financial flows?''

Three results support this reframing. Financial contagion amplifies ROW output losses by 3.45$\times$ relative to the trade-only baseline, and a structural VAR(1) on four historical Middle East conflict episodes confirms that Iran-to-third-country spread transmission of $+0.50$ to $+0.57$ ($t > 5$) is not a calibration artifact. Macroprudential policy that limits spread spillovers would save an order of magnitude more welfare than any plausible monetary adjustment. Varying Taylor-rule aggressiveness from $\phi_\pi = 1.0$ to $\phi_\pi = 2.5$ moves global welfare by less than 0.11 percentage points across nine policy counterfactuals, confirming the impotence of demand-side tools against a 27 percent war-site GDP collapse and a 63 percent oil spike. And within China's stabilization toolkit, the strategic petroleum reserve hits its physical capacity ceiling from the first quarter while manufacturing capacity remains the unconstrained margin---doubling the manufacturing reaction coefficient saves 0.70 percentage points of China's own welfare, whereas additional reserves save nothing at the margin.

Two methodological results strengthen the structural conclusions. The second-order perturbation exercise shows a 5.9 percentage-point linearization bias in Iran's peak loss but leaves the cross-regional spillover pattern intact, so the findings about third-country costs and stabilization rankings are not artifacts of linearization. The VAR-estimated financial linkage matrix moves headline welfare numbers by less than 5 percent relative to the calibrated baseline, confirming that the 4.88 percentage-point financial-channel finding is not sensitive to the specific calibration of $\mathbf{L}$.

Three policy implications follow with precision. The Strait of Hormuz is a critical vulnerability whose costs are dominated by financial transmission, not goods trade---strategic planning that focuses on energy supply continuity without addressing cross-border financial contagion is incomplete. Central banks face a fundamental supply-side constraint in rare-disaster geopolitical scenarios; fiscal coordination, energy diversification, and supply-chain redundancy are the operative policy tools. China's manufacturing capacity generates positive global externalities during geopolitical crises that are larger than the welfare gains from additional petroleum reserves, pointing toward a reinterpretation of China's global economic role that goes beyond the competitive-disruption framing of \citet{autor2013china}.

Future work should extend the framework in three directions: an estimated Bayesian DSGE model to discipline the structural parameters, richer heterogeneous-agent financial frictions to track distributional effects of the contagion, and higher-frequency sovereign CDS data to sharpen the VAR identification of the financial linkage matrix.


% ============================================================================
% REFERENCES
% ============================================================================

\newpage
\bibliographystyle{aer}

\begin{thebibliography}{99}

\bibitem[Abadie and Gardeazabal(2003)]{abadie2003economic}
Abadie, Alberto, and Javier Gardeazabal. 2003. ``The Economic Costs of Conflict: A Case Study of the Basque Country.'' \textit{American Economic Review}, 93(1): 113--132.

\bibitem[Autor, Dorn, and Hanson(2013)]{autor2013china}
Autor, David H., David Dorn, and Gordon H. Hanson. 2013. ``The China Syndrome: Local Labor Market Effects of Import Competition in the United States.'' \textit{American Economic Review}, 103(6): 2121--2168.

\bibitem[Barro(2006)]{barro2006rare}
Barro, Robert J. 2006. ``Rare Disasters and Asset Markets in the Twentieth Century.'' \textit{Quarterly Journal of Economics}, 121(3): 823--866.

\bibitem[Bernanke(2005)]{bernanke2005global}
Bernanke, Ben S. 2005. ``The Global Saving Glut and the U.S.\ Current Account Deficit.'' Sandridge Lecture, Virginia Association of Economists.

\bibitem[Blanchard and Gal\'{i}(2007)]{blanchard2007macroeconomic}
Blanchard, Olivier J., and Jordi Gal\'{i}. 2007. ``The Macroeconomic Effects of Oil Price Shocks: Why Are the 2000s So Different from the 1970s?'' NBER Working Paper 13368.

\bibitem[Bodenstein, Erceg, and Guerrieri(2011)]{bodenstein2011oil}
Bodenstein, Martin, Christopher J. Erceg, and Luca Guerrieri. 2011. ``Oil Shocks and External Adjustment.'' \textit{Journal of International Economics}, 83(2): 168--184.

\bibitem[Borensztein and Panizza(2013)]{borensztein2013debt}
Borensztein, Eduardo, and Ugo Panizza. 2013. ``The Costs of Sovereign Default.'' \textit{IMF Staff Papers}, 56(4): 683--741.

\bibitem[Caldara and Iacoviello(2022)]{caldara2022measuring}
Caldara, Dario, and Matteo Iacoviello. 2022. ``Measuring Geopolitical Risk.'' \textit{American Economic Review}, 112(4): 1194--1225.

\bibitem[{Energy Information Administration}(2023)]{eia2023hormuz}
Energy Information Administration. 2023. ``The Strait of Hormuz Is the World's Most Important Oil Transit Chokepoint.'' EIA Today in Energy, June 20.

\bibitem[Federle et~al.(2026)]{federle2026price}
Federle, Jonathan, Andr\'{e} Meier, Gernot J. M\"{u}ller, Willi Mutschler, and Moritz Schularick. 2026. ``The Price of War.'' \textit{American Economic Review}, 116(3): 791--827.

\bibitem[Gal\'{i}(2015)]{gali2015monetary}
Gal\'{i}, Jordi. 2015. \textit{Monetary Policy, Inflation, and the Business Cycle: An Introduction to the New Keynesian Framework}. 2nd ed. Princeton: Princeton University Press.

\bibitem[Hamilton(1983)]{hamilton1983oil}
Hamilton, James D. 1983. ``Oil and the Macroeconomy since World War II.'' \textit{Journal of Political Economy}, 91(2): 228--248.

\bibitem[Hamilton(2009)]{hamilton2009understanding}
Hamilton, James D. 2009. ``Understanding Crude Oil Prices.'' \textit{Energy Journal}, 30(2): 179--206.

\bibitem[Helveston et~al.(2019)]{helveston2019china}
Helveston, John P., Yimin He, and Michael R. Davidson. 2019. ``Quantifying the Cost Savings of Global Solar Photovoltaic Supply Chains.'' \textit{Nature}, 612: 83--87.

\bibitem[Hilscher and Nosbusch(2010)]{hilscher2010determinants}
Hilscher, Jens, and Yves Nosbusch. 2010. ``Determinants of Sovereign Risk: Macroeconomic Fundamentals and the Pricing of Sovereign Debt.'' \textit{Review of Finance}, 14(2): 235--262.

\bibitem[Jord\`{a}(2005)]{jorda2005estimation}
Jord\`{a}, \`{O}scar. 2005. ``Estimation and Inference of Impulse Responses by Local Projections.'' \textit{American Economic Review}, 95(1): 161--182.

\bibitem[Kilian(2009)]{kilian2009not}
Kilian, Lutz. 2009. ``Not All Oil Price Shocks Are Alike: Disentangling Demand and Supply Shocks in the Crude Oil Market.'' \textit{American Economic Review}, 99(3): 1053--1069.

\bibitem[Longstaff et~al.(2011)]{longstaff2011financial}
Longstaff, Francis A., Jun Pan, Lasse H. Pedersen, and Kenneth J. Singleton. 2011. ``How Sovereign Is Sovereign Credit Risk?'' \textit{American Economic Journal: Macroeconomics}, 3(2): 75--103.

\bibitem[Lucas(1987)]{lucas1987models}
Lucas, Robert E., Jr. 1987. \textit{Models of Business Cycles}. Oxford: Basil Blackwell.

\bibitem[Ramey(2011)]{ramey2011identifying}
Ramey, Valerie A. 2011. ``Identifying Government Spending Shocks: It's All in the Timing.'' \textit{Quarterly Journal of Economics}, 126(1): 1--50.

\bibitem[Rey(2015)]{rey2015dilemma}
Rey, H\'{e}l\`{e}ne. 2015. ``Dilemma Not Trilemma: The Global Financial Cycle and Monetary Policy Independence.'' NBER Working Paper 21162.

\bibitem[Zhang et~al.(2020)]{zhang2020china}
Zhang, Dayong, Qiang Ji, and Apostolos Kiohos. 2020. ``China's Oil Strategic Petroleum Reserve: A DSGE Analysis.'' \textit{Energy Economics}, 88: 104748.

\end{thebibliography}


% ============================================================================
% APPENDIX (FOR ONLINE PUBLICATION)
% ============================================================================

\newpage
\appendix
\section*{Appendix: Model Equations Summary}
\addcontentsline{toc}{section}{Appendix}

For reference, we collect the complete system of equations for each region $i$:

\begin{align}
\pi_{i,t} &= \beta \, \mathbb{E}_t[\pi_{i,t+1}] + \kappa_i \, \hat{y}_{i,t-1} + \xi^{\text{supply}}_{i,t} \tag{A.1} \\[6pt]
\hat{y}_{i,t} &= \rho \, \hat{y}_{i,t-1} - \frac{1}{\sigma_i}(i_{i,t} - \mathbb{E}_t[\pi_{i,t+1}]) + 0.8 \, g^{\text{mil}}_{i,t} - \tau^{\text{trade}}_{i,t} - 0.3 \, \alpha_{e,i} \Delta p^{\text{oil}}_t + \Phi^{\text{fin}}_{i,t} \tag{A.2} \\[6pt]
i_{i,t} &= \rho_{i,i} \, i_{i,t-1} + (1-\rho_{i,i})(\phi_{\pi,i} \pi_{i,t} + \phi_{y,i} \hat{y}_{i,t-1}) \tag{A.3} \\[6pt]
s_{i,t} &= s^{\text{shock}}_{i,t} + 0.3 \max(-\hat{y}_{i,t-1}, 0) + 0.2 \, s_{i,t-1} \tag{A.4} \\[6pt]
k_{i,t} &= k^{\text{shock}}_{i,t} - 2 s_{i,t} + 0.5(i_{i,t} - 0.02) + 0.3 \, k_{i,t-1} \tag{A.5} \\[6pt]
q_{i,t} &= 0.6 \, q_{i,t-1} + 10 g^e_{i,t} - 8(i_{i,t} + s_{i,t}) - 3 \max(-\hat{y}_{i,t},0) + q^{\text{shock}}_{i,t} \tag{A.6} \\[6pt]
e_{i,t} &= 0.4 \, e_{i,t-1} + 0.5 \, s_{i,t} - 0.3 \, k_{i,t} + \iota^{\text{fx}}_{i,t} \tag{A.7} \\[6pt]
\Phi^{\text{fin}}_{i,t} &= -0.15 \, s_{i,t} - 0.10 \max(-k_{i,t}, 0) \tag{A.8}
\end{align}

The global oil market clears according to:
\begin{align}
\Delta S_t &= \lambda^{\text{block}}_t (s^H - s^{\text{Iran}}) + s^{\text{Iran}} \lambda^{\text{prod}} - 0.3 \lambda^{\text{block}}_t (s^H - s^{\text{Iran}}) \tag{A.9} \\[6pt]
\frac{\Delta P^{\text{oil}}_t}{P^{\text{oil}}} &= \frac{\Delta S_t - \text{SPR}_t + \Delta D_t}{|\varepsilon_d| + \varepsilon_s} \tag{A.10}
\end{align}

Financial contagion operates through:
\begin{equation}
\tilde{s}^{\text{shock}}_{j,t} = s^{\text{shock}}_{j,t} + 0.1 \Big(\textstyle\sum_{i \neq j} L_{ij} s_{i,t-1}\Big)(1 + 2\max(-k_{\text{Iran},t}, 0)) \tag{A.11}
\end{equation}

\end{document}
